\subsection{Directed and Undirected Graphical Models}

Graphical models use graph structure to encode conditional independence and to organize high-dimensional probability laws.

\begin{definition} \label{def:sl-directed-graphical-model} \hlt{Directed Graphical Model}\\
For a directed acyclic graph $G$ with vertices $1,\ldots,d$, a distribution factorizes according to $G$ if
\[
p(x_1,\ldots,x_d)=\prod_{j=1}^d p(x_j\mid x_{\operatorname{pa}(j)}),
\]
where $\operatorname{pa}(j)$ denotes the parents of vertex $j$.
\end{definition}

\begin{definition} \label{def:sl-markov-random-field} \hlt{Undirected Graphical Model}\\
For an undirected graph $G$, a positive distribution is Markov with respect to $G$ when missing edges encode conditional independence.  It often factorizes over cliques as
\[
p(x)=\frac1Z\prod_{C}\psi_C(x_C),
\]
where $Z$ is the normalizing constant.
\end{definition}

\begin{remark} \label{rem:sl-separation-graphical-models} \hlt{Separation}\\
Graph separation rules translate paths in the graph into conditional independence statements.  Directed models use d-separation; undirected models use ordinary graph separation.
\end{remark}
